% !PRINTED 366
\weberpage{366}
\gesperrt{If $E'$ is the group of the units contained in $O'$, and if $(E, E')$ has a finite value, then the class number $h'$ is finite and determined by the formula (12).}

If, for example, $O$ is the group of all numbers in $\Omega$ (apart from 0), and $O'$ the group of the numbers with positive norm, then, provided there are any numbers with negative norm at all, $(O, O') = 2$. If there are then units with negative norm, $(E, E') = 2$ as well. But if all units have positive norm, then $(E, E') = 1$, and we obtain $h' = h$ in the first case, $2h' = h$ in the second case.

\gesperrt{The equivalence with respect to $O$ is then the general equivalence, that with respect to $O'$ the sharpened one.}

Let us return to the orders $[Q]$ and understand by $O$ the numbers coprime to $Q$ (if necessary only those with positive norm); then we have determined $(O, O')$ in the preceding section.

If the discriminant $D$ and its stem $\Delta$ are negative, then in general, both in $E$ and in $E'$, there are only the two units $\pm 1$. In the two exceptional cases $\Delta = -3, -4$, $E$ contains six and four units, but $E'$, if $Q > 1$, only two, and thus $(E, E') = 1$ in general, $(E, E') = 3$ in the case $\Delta = -3$, $(E, E') = 2$ in the case $\Delta = -4$. Hence in these cases we have, by \S\,98, (10):
\weq{15}{\lambda h' = Q\varPi\left(1 - \frac{(\Delta, q)}{q}\right)h,}
in which:
\weq{16}{
\begin{aligned}
\lambda &= 2, & &\text{for } \Delta = -4 \\
\lambda &= 3, & &\text{\glqq} \ \Delta = -3 \\
\lambda &= 1, & &\text{\glqq} \ \Delta < -4
\end{aligned}
}
This is the relation which was first derived by \gesperrt{Gauss} and later confirmed by \gesperrt{Dirichlet} in another way.

If, for example, $\Delta \equiv 1 \, (\mathrm{mod}\ 4)$, $Q = 2$, then in \gesperrt{Gaussian} notation $h$ is the class number for the forms of the second kind, $h'$ that for the forms of the first kind. We have
\[
\begin{aligned}
Q\left(1 - \frac{(\Delta, q)}{2}\right) &= 1, & \Delta &\equiv 1 \, (\mathrm{mod}\ 8), \\
&= 3, & \Delta &\equiv 5 \, (\mathrm{mod}\ 8).
\end{aligned}
\]
Thus
\[
\begin{aligned}
h' &= h, & \Delta &\equiv 1 \\
&= 3h, & \Delta &\equiv 5
\end{aligned}
\qquad (\mathrm{mod}\ 8),
\]
and only in the exceptional case $\Delta = -3$ does $h' = h$ occur.
